% GENERATED FILE. Edit canonical lessons, then run npm run generate:books.
% canonical-source-hash: fnv1a64:2859ef2405714ff2
% canonical-lessons: JA-C364-songai, JA-C364-benshou-suru, JA-C364-osameru, JA-C364-ukemotsu, JA-C364-sekinin

\chapter{Damage and Compensation}
\label{ch:ja-damage-and-compensation}
\hlchaptermodality{364}

\begin{chapteropening}
\textbf{By the end of this chapter:} \emph{I can say who pays if something is damaged.}

Complete the last lesson of chapter 364: I can say who pays if something is damaged.
\end{chapteropening}

\section[\texorpdfstring{\ja{そんがい} (songai)}{そんがい (songai)}]{\ja{そんがい} (songai) — damage; a loss}

\label{lesson:JA-C364-songai}



\begin{quote}
Before the new one: say the Japanese for half price, then the Japanese for luxury.
\end{quote}

\subsection*{You'll want to know: \ja{そんがい}}
\textbf{\ja{そんがい}} — \emph{songai} — \textquotedblleft{}damage; a loss\textquotedblright{}.

\begin{quote}
\textbf{\ja{たいふうが} \ja{きたら、そんがいが} \ja{あるでしょう。}} — \emph{taifū ga kitara, songai ga aru deshō} — \textquotedblleft{}If the typhoon comes, there will probably be damage.\textquotedblright{}
\end{quote}

\begin{cousinweb}
\textbf{\ja{そんがい}} (\emph{songai}) is Sino-Japanese (kango): two Chinese characters read with their on'yomi, readings borrowed from Chinese. They are \emph{son}, \textquotedblleft{}loss\textquotedblright{} (Middle Chinese \emph{swonX}) and \emph{gai}, \textquotedblleft{}harm\textquotedblright{} (\emph{hajH}) (KANJIDIC2; Middle Chinese in Baxter's transcription).
\end{cousinweb}

\begin{grammarlens}[title={What you've built}]
One more word for saying what happens if.
\end{grammarlens}

\subsection*{Your turn}
\begin{itemize}
\raggedright
  \item \emph{Say it:} \emph{songai}
  \item \emph{Say it:} \emph{songai}, once more
  \item \emph{Say it:} \emph{taifū ga kitara, songai ga aru deshō} — If the typhoon comes, there will probably be damage.
  \item \emph{Recall:} say the Japanese for half price, then the Japanese for luxury, then say \emph{songai} again
\end{itemize}

\begin{tcolorbox}[breakable,colback=teal!4,colframe=teal!35!black,title={Before you move on}]
What does \ja{そんがい} mean? (\textquotedblleft{}Damage; a loss\textquotedblright{}.) Say it once more.
\end{tcolorbox}
\section[\texorpdfstring{\ja{べんしょうする} (benshō suru)}{べんしょうする (benshō suru)}]{\ja{べんしょうする} (benshō suru) — to pay for damage; to compensate}

\label{lesson:JA-C364-benshou-suru}



\begin{quote}
Before the new one: say the Japanese for luxury, then the Japanese for damage to property.
\end{quote}

\subsection*{You'll want to know: \ja{べんしょうする}}
\textbf{\ja{べんしょうする}} — \emph{benshō suru} — \textquotedblleft{}to pay for damage; to compensate\textquotedblright{}.

\begin{quote}
\textbf{\ja{おさらを} \ja{わったら、べんしょうして} \ja{ください。}} — \emph{osara o wattara, benshō shite kudasai} — \textquotedblleft{}If you break a plate, please pay for it.\textquotedblright{}
\end{quote}

Polite \textbf{\ja{べんしょうします}}, \emph{benshōshimasu}; a \ja{する} verb. Past (polite): \textbf{\ja{きゃくが} \ja{べんしょうしました。}}, \emph{kyaku ga benshō shimashita}.

\begin{cousinweb}
\textbf{\ja{べんしょうする}} (\emph{benshō suru}) is built from a Sino-Japanese noun and the native verb \emph{suru}, \textquotedblleft{}to do\textquotedblright{}. The noun, \emph{benshō}, is two Chinese characters read with their on'yomi, readings borrowed from Chinese: \emph{ben}, \textquotedblleft{}distinguish\textquotedblright{} and \emph{shō}, \textquotedblleft{}reparation\textquotedblright{} (KANJIDIC2; JMdict).
\end{cousinweb}

\begin{grammarlens}[title={What you've built}]
One more word for saying what happens if.
\end{grammarlens}

\subsection*{Your turn}
\begin{itemize}
\raggedright
  \item \emph{Say it:} \emph{benshō suru}
  \item \emph{Say it:} \emph{benshō suru}, once more
  \item \emph{Say it:} \emph{osara o wattara, benshō shite kudasai} — If you break a plate, please pay for it.
  \item \emph{Recall:} say the Japanese for luxury, then the Japanese for damage to property, then say \emph{benshō suru} again
\end{itemize}

\begin{tcolorbox}[breakable,colback=teal!4,colframe=teal!35!black,title={Before you move on}]
What does \ja{べんしょうする} mean? (\textquotedblleft{}To pay for damage; to compensate\textquotedblright{}.) Say it once more.
\end{tcolorbox}
\section[\texorpdfstring{\ja{おさめる} (osameru)}{おさめる (osameru)}]{\ja{おさめる} (osameru) — to pay taxes; to pay (a fee); to put away}

\label{lesson:JA-C364-osameru}



\begin{quote}
Before the new one: say the Japanese for damage to property, then the Japanese for to pay for damage.
\end{quote}

\subsection*{You'll want to know: \ja{おさめる}}
\textbf{\ja{おさめる}} — \emph{osameru} — \textquotedblleft{}to pay taxes; to pay (a fee); to put away\textquotedblright{}.

\begin{quote}
\textbf{\ja{らいげつに} \ja{なったら、ぜいきんを} \ja{おさめます。}} — \emph{raigetsu ni nattara, zeikin o osamemasu} — \textquotedblleft{}Next month, I will pay my taxes.\textquotedblright{}
\end{quote}

Polite \textbf{\ja{おさめます}}, \emph{osamemasu}; an ichidan verb. Past (polite): \textbf{\ja{ちちが} \ja{ぜいきんを} \ja{おさめました。}}, \emph{chichi ga zeikin o osamemashita}.

\begin{cousinweb}
\textbf{\ja{おさめる}} (\emph{osameru}) is a native Japanese word (wago), inherited rather than borrowed from Chinese, and its older origin is not settled. Its kanji was attached to it later as a kun'yomi, a native reading given to a Chinese character (JMdict, KANJIDIC2).
\end{cousinweb}

\begin{grammarlens}[title={What you've built}]
One more word for saying what happens if.
\end{grammarlens}

\subsection*{Your turn}
\begin{itemize}
\raggedright
  \item \emph{Say it:} \emph{osameru}
  \item \emph{Say it:} \emph{osameru}, once more
  \item \emph{Say it:} \emph{raigetsu ni nattara, zeikin o osamemasu} — Next month, I will pay my taxes.
  \item \emph{Recall:} say the Japanese for damage to property, then the Japanese for to pay for damage, then say \emph{osameru} again
\end{itemize}

\begin{tcolorbox}[breakable,colback=teal!4,colframe=teal!35!black,title={Before you move on}]
What does \ja{おさめる} mean? (\textquotedblleft{}To pay taxes; to pay (a fee); to put away\textquotedblright{}.) Say it once more.
\end{tcolorbox}
\section[\texorpdfstring{\ja{うけもつ} (ukemotsu)}{うけもつ (ukemotsu)}]{\ja{うけもつ} (ukemotsu) — to be in charge of}

\label{lesson:JA-C364-ukemotsu}



\begin{quote}
Before the new one: say the Japanese for to pay for damage, then the Japanese for to pay taxes.
\end{quote}

\subsection*{You'll want to know: \ja{うけもつ}}
\textbf{\ja{うけもつ}} — \emph{ukemotsu} — \textquotedblleft{}to be in charge of\textquotedblright{}.

\begin{quote}
\textbf{\ja{せんぱいが} \ja{やめたら、あにが} \ja{しごとを} \ja{うけもつでしょう。}} — \emph{senpai ga yametara, ani ga shigoto o ukemotsu deshō} — \textquotedblleft{}If my senior leaves, my brother will probably take charge of the work.\textquotedblright{}
\end{quote}

Polite \textbf{\ja{うけもちます}}, \emph{ukemochimasu}; a godan verb. Past (polite): \textbf{\ja{あにが} \ja{しごとを} \ja{うけもちました。}}, \emph{ani ga shigoto o ukemochimashita}.

\begin{cousinweb}
\textbf{\ja{うけもつ}} (\emph{ukemotsu}) is a native Japanese word (wago), inherited rather than borrowed from Chinese: it is built from \emph{uke}, \textquotedblleft{}receiving (ukeru)\textquotedblright{} and \emph{motsu}, \textquotedblleft{}to hold\textquotedblright{}, and its kanji were attached later as kun'yomi (JMdict).
\end{cousinweb}

\begin{grammarlens}[title={What you've built}]
One more word for saying what happens if.
\end{grammarlens}

\subsection*{Your turn}
\begin{itemize}
\raggedright
  \item \emph{Say it:} \emph{ukemotsu}
  \item \emph{Say it:} \emph{ukemotsu}, once more
  \item \emph{Say it:} \emph{senpai ga yametara, ani ga shigoto o ukemotsu deshō} — If my senior leaves, my brother will probably take charge of the work.
  \item \emph{Recall:} say the Japanese for to pay for damage, then the Japanese for to pay taxes, then say \emph{ukemotsu} again
\end{itemize}

\begin{tcolorbox}[breakable,colback=teal!4,colframe=teal!35!black,title={Before you move on}]
What does \ja{うけもつ} mean? (\textquotedblleft{}To be in charge of\textquotedblright{}.) Say it once more.
\end{tcolorbox}
\section[\texorpdfstring{\ja{せきにん} (sekinin)}{せきにん (sekinin)}]{\ja{せきにん} (sekinin) — responsibility}

\label{lesson:JA-C364-sekinin}



\begin{quote}
Before the new one: say the Japanese for to pay taxes, then the Japanese for to be in charge of.
\end{quote}

\subsection*{You'll want to know: \ja{せきにん}}
\textbf{\ja{せきにん}} — \emph{sekinin} — \textquotedblleft{}responsibility\textquotedblright{}.

\begin{quote}
\textbf{\ja{もし} \ja{こわしたら、せきにんを} \ja{とって} \ja{ください。}} — \emph{moshi kowashitara, sekinin o totte kudasai} — \textquotedblleft{}If you break it, please take responsibility.\textquotedblright{}
\end{quote}

\begin{cousinweb}
\textbf{\ja{せきにん}} (\emph{sekinin}) is Sino-Japanese (kango): two Chinese characters read with their on'yomi, readings borrowed from Chinese. They are \emph{seki}, \textquotedblleft{}blame\textquotedblright{} (Middle Chinese \emph{tsreak}) and \emph{nin}, \textquotedblleft{}duty\textquotedblright{} (\emph{nyim}) (KANJIDIC2; Middle Chinese in Baxter's transcription).
\end{cousinweb}

\begin{grammarlens}[title={What you've built}]
One more word for saying what happens if.
\end{grammarlens}

\subsection*{Your turn}
\begin{itemize}
\raggedright
  \item \emph{Say it:} \emph{sekinin}
  \item \emph{Say it:} \emph{sekinin}, once more
  \item \emph{Say it:} \emph{moshi kowashitara, sekinin o totte kudasai} — If you break it, please take responsibility.
  \item \emph{Recall:} say the Japanese for to pay taxes, then the Japanese for to be in charge of, then say \emph{sekinin} again
\end{itemize}

\begin{tcolorbox}[breakable,colback=teal!4,colframe=teal!35!black,title={Before you move on}]
What does \ja{せきにん} mean? (\textquotedblleft{}Responsibility\textquotedblright{}.) Say it once more.
\end{tcolorbox}
